% !PRINTED 414
\weberpage{414}
Ist dagegen die Gleichung (1) nicht eine identische, so ist $\omega$
die Wurzel einer ganzzahligen quadratischen Gleichung:
\weq{2}{b\,\omega^2 + (a-\partial)\,\omega - c = 0,}
also, wenn wir
\[
a\,\partial - b\,c = n
\]
\[
m = -4\,b\,c - (a-\partial)^2 = 4\,n - (a+\partial)^2
\]
setzen, so ist
\[
m \equiv 0 \ \text{ oder } \ \equiv -1 \ (\mathrm{mod}\ 4),
\]
und es ergibt sich:
\[
\omega = \frac{-a+\partial+\sqrt{-m}}{2\,b},
\]
\[
\mu = \frac{a+\partial+\sqrt{-m}}{2},
\]
\[
n = \frac{a+\partial+\sqrt{-m}}{2}\cdot\frac{a+\partial-\sqrt{-m}}{2}, % Source scan shows a minus; the mark above it is a scan/print speck.
\]
\[
\mu^2 - (a+\partial)\,\mu + n = 0.
\]

Damit $\omega$ einen positiven imaginären Teil habe, muß $m$ und
umsomehr also $n$ \gesperrt{positiv} sein, und $\mu$ ist eine komplexe, ganze,
algebraische Zahl. Daher erklärt sich die Bezeichnung \gesperrt{komplexe
Multiplikation}.

Wenn nun umgekehrt $\omega$ einer quadratischen Gleichung
\weq{3}{A\,\omega^2 + B\,\omega + C = 0}
genügt, worin $A$, $B$, $C$ ganze Zahlen sind, und die Diskriminante
\weq{4}{D = B^2 - 4\,A\,C}
negativ ist, so ist $\omega$ nicht reell, und in einer der beiden Wurzeln
von (3) ist der imaginäre Teil positiv. Diese soll für $\omega$ genommen
werden. Es heißt dann $\omega$ eine Wurzel der quadratischen Form
$(A, B, C)$.

Unbeschadet der Allgemeinheit können $A$, $B$, $C$ ohne gemeinsamen
Teiler und $A$, $C$ positiv angenommen werden. Es lassen sich dann für
$\omega$ unendlich viele Relationen von der Form (1) [oder (2)]
aufstellen. Man hat nur, wenn $x$ eine ganze von Null verschiedene Zahl
ist, zu setzen:
\weq{5}{\begin{aligned}
b &= A\,x,\\
c &= -\,C\,x,\\
a-\partial &= B\,x.
\end{aligned}}
